\section{The Cosmic Revelation}

\begin{quotex}
Without Christianity I don't think the oriental religions, Hinduism and Buddhism, can answer the needs of the modern world. But without the enrichment of the mystical tradition of Asia I doubt whether the Western Churches can really discover the fullness of Christ which we are seeking. \flright{\textsc{Bede Griffiths}}

\end{quotex}
Bede Griffiths, a one-time student of C S Lewis, combined the life of a monk and a yogi. On his own, apparently, he undertook Guenon's project to reinvigorate his Catholic tradition with Oriental metaphysics. In a series of six lectures, published as \emph{The Cosmic Revelation}, Griffiths provides us with a fine overview of the Vedanta and its relevance to the West. I'll organize the following review by the chapter titles while highlighting their main points.

\paragraph{The Vedic Revelation}
As does \textbf{Rene Guenon}, Griffiths accepts that the \emph{Vedas} arise from a divine revelation. In them, the word (\emph{vac}) is the mediator between God and man. The word \emph{veda} means knowledge, the same as the Latin \emph{video}, to see, to know. Hence ``knowing" is like ``seeing". As Griffiths points out, the Vedas (and Sanskrit) were brought to India by the Aryans from the North in the second millennium BC. Its authors are called rishis, or see-ers. But Vedas are not a fabrication; they are also called \emph{sruti}, which means they were ``heard". In sum, they are what has been heard, what is eternal, without human authorship.

After the Vedas came the \emph{Brahmanas}, which were commentaries on the ritual sacrifices, the \emph{Aranyakas}, written by the mystical forest dwellers, and then the Upanishads, or records of discussions between the rishis and their students. These are all the echoes of the original peoples who arrived in India from the North; that is why they may still resonate in some deep part of us.

The later epics (the \emph{Mahabharata} and the \emph{Ramayana}) evolved over time from the intermixtures of the Northern Aryans with the Southern Dravidians. Griffiths notes that the northern religion is ``a religion of a nomadic people, of warriors, and much more dynamic", in contrast with the earthy religion of the Mother goddess worship of the South. Next came the \emph{Darshanas}, or the orthodox philosophical systems.

\subparagraph{Monotheism}
Hinduism, he says, is called the \emph{sanatana dharma}, the ``eternal religion". Griffiths does not mention that St Augustine made the same point: there is but one tradition, which is now known as Catholic. Therefore, Griffiths needs to explain the two concepts of God. The first Veda refers to ``the one being the wise called by many names." Therefore, the gods are names and forms under which the one God manifests himself. That is, they are more or less the equivalent of ``angels" or the ``cosmic powers" mentioned by St Paul. The One Being is behind the powers of nature, the sky, the earth, the sea, fire, and so on. Hence, Hinduism is no more polytheistic than is Catholicism with its cult of angels and saints.

\subparagraph{Esoteric Understanding.}
Words in sacred languages are not as univocal as modern ones. For example, the word ``spirit" in most languages also means ``wind or air or breath or life or soul or spirit." That is what makes translations difficult: there is no one precise meaning, since the word contains all the meanings. Griffiths refers to \emph{Poetic Diction} by the anthroposophist and inkling \textbf{Owen Barfield} for support.

The rishis did not have the same dualistic consciousness of most people today. For example, in the Vedas we find many prayers for material prosperity: cattle, horses, butter, wealth, and so on. \textbf{Aurobindo Ghose}, in the \textit{Secret of the Veda}\footnote{\url{https://www.sriaurobindoashram.org/sriaurobindo/writings.php}} points out that these words have both a psychological and a physical meaning. So as butter comes from milk and thoughts come from mind, ``butter" may also mean thought.

A better example, from our perspective, is the fire sacrifice associated with the hearth, since keeping the ``home fire burning" was such an important rite to early Indo-European people. But that also means to build the fire in the heart. \emph{Tapas}, which means ``asceticism, self-control, discipline", originally meant ``heat". Hence, the buildup of friction in the mind through self-discipline is the same as ``heart" or the fire in the heart. So there is the exoteric religion of the rites around the physical fire, and the esoteric religion of inner fire. Yet it would be a mistake to separate them, or to claim that the esoteric meaning is the ``real" meaning. That is a dualistic view. Rather, our ancestors would have experienced them as the same; the fire in the hearth and the fire in the heart are the same. The discipline of the outer rite is a reflection of the inner discipline.

\subparagraph{The Three Worlds}
The rishis recognized the three worlds of physical being, psychological being, and spiritual being as one, not separating them, as we may do. (Actually, the spiritual is now opaque to most people, and even the psychological is understood as a mere physical process.) Hence, the sun is not just an astronomical body providing light to the physical world but also the source of light to the mind as the inner sun. Hence, the sun, moon, and earth, trees, etc., understood at the level of soul and spirit, are revered because God is in the midst of them. That is, the whole creation is pervaded by God.

\subparagraph{Two Conceptions of God}
The Hebrews begin with God's transcendence. God is infinitely high above the heavens and descends to the prophets, becomes incarnate in Jesus, and so on. That is, it is a movement of descent from above.

The Hindu starts in the opposite way. God is immanent in every created thing. Of course, Christians pay lip service to the idea of omnipresence, but have largely lost the experience of it. Hence, these two visions are complementary, not opposed.

Of course, God as transcendent to the physical cosmos is one thing, but God as transcendent spiritually is another. Hence, we should also understand transcendence in our interiority, in our ``heart", as the I that watches the flow of consciousness. Many theological conundrums can be resolved through spiritual experience and understanding.

Griffiths concludes the first lecture with a call for the Church to be enriched from the encounter with other religions, just as it absorbed European paganisms. However, I believe he is neglecting a certain fundamental point. Europeans and the people who were formed around the Vedic revelation have the same ancestors. Hence, there is a certain affinity of the Vedic revelation to the West which makes it more appropriate.

As the Church has been expanding into Africa and elsewhere, there is an encounter, as he points out, with those aboriginal religions, leading, for example, to cults like Santeria. However, as much as it may be valid for them, that is not our present concern.

\paragraph{The Cosmic Covenant}

Griffiths next tries to relate the Vedic Revelation to Biblical Revelation. Biblical history is punctuated with a series of covenants. The first one was God's covenant with Adam, which just means ``man". Since God was revealed to all mankind, Griffiths calls it the Cosmic Covenant. Its origin is in Paradise, which is that state where man is in harmony with nature and God.

This harmony has been broken. For example, ecological issues reflect the disharmony of man with nature. Wars, crime, conflict, etc., reveal the lack of harmony of man with man. Finally, the communion with God has been lost. The loss of the primordial harmony is called Original Sin. Salvation, then, is the restoration of that primordial harmony.

The next covenant was with Noah, from whose descendants all the nations of the world arose. As the father of mankind, this was a re-expression of the Cosmic Covenant. Thus, all the valid Traditions are rooted in that covenant. There are righteous pagans in the Bible, such as Enoch.

Griffiths then turns to Melchizedek, the pagan priest and king of Salem. In his conversation with Abraham, Melchizedek appeals to the ``most high God, the creator of heaven and earth." Griffiths takes this to mean that there was a revelation to the pagans, specifically mentioning the Amerindians and the Hindus. He makes the interesting point that Christ was a priest in the order of Melchizedek, not in the order of Hebrew priests starting with Aaron.

Griffiths provides scriptural support for this claim. In Paul's letter to the Romans, we read that there are two sources of revelation to the pagans.

\begin{enumerate}
\item The creation of the world. Paul writes, ``His invisible nature is made known by the visible things He has made." (Romans 1:20) 
\item Within the Heart of man. Again, Paul writes, ``The Pagans who have not the law of Moses have a law written in their hearts: their conscience accusing or perhaps excusing them until the day when God will judge the world by Jesus Christ." (Romans 2:14) 
\end{enumerate}
In other words, these are the revelations of the macrocosm and the microcosm, and are relatable to the Pythagorean and Hermetic initiations described by Valentin Tomberg in Letter VI of \emph{Meditations on the Tarot}.

\paragraph{The Sky Father}
Early man had not acquired the habit of abstract thought, as his mode of thinking was more intuitive. As was pointed out earlier, he saw the physical, the psychological, and the spiritual worlds in their integrated wholeness. \textbf{Mircea Eliade}, in \emph{Patterns of Comparative Religion}, shows how primitive man could come to an almost total understanding of God simply by looking at the sky. Griffiths summarizes it this way:

\begin{quotex}
Consequently, when he [early man] looked up at the sky, the sky was a revelation to him. It was not merely the physical phenomenon which we see; to him the sky had a psychological meaning and a spiritual meaning. It revealed the creator present in his creation. 

\end{quotex}
Therefore, there are several things to be intuited.

\begin{itemize}
\item The sky is vast and infinitely above us. This is a revelation of God. 
\item The sky has no apparent limit, unlike all other objects we experience. The sky is all-embracing, without limit, revealing the idea of infinity. 
\item Everything in the world is changing. Life arises and dies, the sun rises and sets, the seasons change, etc. Yet the sky is unchanging, just like God. 
\item There is an order in the sky, so by analogy there is a Cosmic Order. 
\end{itemize}
Griffiths gives examples of how the image of God as immense, infinite, eternal, the source of order, providence and fear, arises, based on the experience of nature (e.g., rain $\rightarrow$ providence, thunder $\rightarrow$ fear). We have discussed several times the method of Hermetic meditation in order to recover that primordial state. Also, learn to use the analogy of being to see the connectedness between the physical, psychological, and spiritual worlds.

The Vedas mention the ancient God ``\emph{Dyaus Pita}", or literally ``Sky Father". In Greek that becomes ``\emph{Zeus Pater}", then in Latin, ``\emph{Dios Pater}", or ``Jupiter". Hence we see, through philology, how the Hyperborean migrations related to each other. So, when Jesus starts his prayer with ``Our Father in the Sky", we see a reference to that most ancient revelation. (``Heaven" simply means ``sky"; this is clearer in other languages.)

\paragraph{Angels and Cosmic Powers}
Although our world conception has been devitalized in the modern world through the materialization of our ideas of heaven and earth, as late as the Middle Ages, men believed that God ruled the world through the angels, or ``cosmic powers". The angels were present in all creation, especially the stars, sun, moon, and the ``host of heaven." ``The sky was not simply a material thing, it was the abode of the gods, the manifestation of the Supreme Spirit, the Supreme Father, who ruled over all." Griffiths emphasizes that we have not ``progressed" beyond that understanding; rather, the modern view is actually a devolution. He then uses the following example from Hinduism, the highest expression of the cosmic revelation.

\begin{quotex}
There we have the worship of God through the power of nature. The devas, literally ``the shining ones," are these cosmic powers manifesting themselves in the whole creation. In Hinduism the supreme Reality is known as Brahman, and this Brahman, the One Being, manifests himself in the devas, the gods, the powers which rule the creation. 

\end{quotex}
Hence, there is the hierarchy, or chain, of Being just as the Medievals believed. This hierarchy starts with the Supreme Being who transcends all, then the devas, the powers of nature, man with intelligence and sense, and finally the world below him. The difference is that this understanding has been lost in the West.

\paragraph{Sacrifice}
Griffiths next goes into some detail about the Hindu mode of worship. Apparently his order developed what he calls the ``Indian rite" of the mass. First, water is used to symbolize purification. Then there is the bread and wine, which represent the fruits of the earth and the work of human hands. Flowers are placed around them to signify that the sacrifice is at the centre of the universe. Then they incense the gifts and wave fire as the flame of burning camphor. Christ, then, is the sacrifice of the whole creation to the Father.

\paragraph{The Presence of God}
Griffiths offers a helpful interpretation of the Hindu trinity of Brahma, Vishnu, and Shiva, as three forms of God. Vishnu as preserver pervades the whole universe. \textbf{Thomas Aquinas} describes how God is in creation this way:

\begin{itemize}
\item ``God is in the whole creation by His power, because He sustains everything by his power." The Hindus call this power ``Shakti". 
\item Since there is no ``spatial" component in God, He does not exercise his power at a distance. Hence, he is in everything by His Presence. But not in the sense that everything is a ``part" of God, who has no parts. Rather, He is present in His essence. 
\end{itemize}
Once again, we have lost this sense of the presence of God, as totally pervaded by God. Several years ago, a woman told me she envisaged God as like a ``gas"; of course, that is simplistic yet I was struck at the time that it indicated some level of understanding of that pervading presence.

\paragraph{The Triple Character of God}
Griffiths lists these as the triple character in God:

\begin{itemize}
\item \textbf{Nirguna Brahman}. That is without qualities, without attributes, beyond everything which can be conceived. 
\item \textbf{Saguna Brahman}. God with attributes, conceived as Creator, Lord, Saviour. 
\item \textbf{Atman}. This is God dwelling in each person as his own inner spirit. 
\end{itemize}
In Christian terms, the first is the Father, the Source, the Beyond (Nirguna Brahman). That Father manifests Himself in the Son, or the Word, the self-manifestation of God (Saguna Brahman). Finally God is present in man's innermost spirit, the Spirit of God (Atman).

Griffiths is not trying to convert anyone to Hinduism, or a crypto-Hinduism, but rather he uses it as a tool to recover our own past.

\paragraph{The Cosmic Mystery}

\begin{quotex}
That Self, which is free from sin, free from old age, from death and grief, from hunger and thirst, which desires nothing but what it ought to desire, and imagines nothing but what it ought to imagine, it is that which we should search out; that which we must try to understand. He who reaches that Self and understands it gains all the world of desires.\flright{\emph{Chandogya Upanishad}}

\end{quotex}
\emph{Lead the thoughts from the head into the heart and keep them there.}

With this saying of the Greek fathers, \textbf{Bede Griffiths} describes the opening of the Heart to the Cosmic Mystery. The Upanishads likewise need to be read that way. Although there are 108 Upanishads in theory, the twelve principle Upanishads are what matter.

From the \emph{Brihadaranyaka Upanishad}, Griffiths identifies three words that signify the Godhead: Brahman, Atman, Purusha.

\begin{itemize}
\item \textbf{Brahman}. Brahman is the source of all creation or the ground of all being. 
\item \textbf{Atman}. Atman, the Self, is the ground of all consciousness. 
\item \textbf{Purusha}. God as Person. 
\end{itemize}
There is a progressive realization that Brahman, Atman, and Purusha are all One. This chapter is concerned with the first two. The understanding of the word ``Brahman" illustrates Griffiths' powerful methodology.

The root of Brahman is ``Brh" which means ``to grow or to swell", which is the physical meaning. The psychological and spiritual meaning then identifies this ``swelling" as the rising up of the awareness of God. Hence, he explains:

\begin{quotex}
The word Brahman came to mean a prayer, something which rises up in the heart, swells within, breaks out and opens up to the divine, to a mystery beyond. Man is in search of this hidden mystery, and it is a mystery which cannot be named. 

\end{quotex}
\paragraph{Sacrifice}
Griffiths reminds us that ``sacrifice was the centre of all ancient religion." Since everything comes down from above, it must be returned. Thus, the sacrifice is the return to God, and sin is the opposite, the appropriation of something to one's self. This rhythm of the world — coming forth from, and then returning to, God — is \emph{rita}, or cosmic order. When you live in that rhythm, you are said to be turning the wheel of the law, the dharma chakra. When you sin, you are going against the law of the universe.

Brahman is the power which sustains the sacrifice and sustains the whole creation. There are two aspects.

\begin{enumerate}
\item Brahman is the source of everything. 
\item Brahman pervades the universe. 
\end{enumerate}
Griffiths rejects the misunderstanding that this represents a form of pantheistic monism. Griffiths then discusses several stories from the Vedas, which we can't summarize here, other than to point out that some truths can only be expressed by stories, parables, myths, and the like.

\paragraph{Positive and Negative Approaches}
Since Brahman can never be fully encompassed in words, two methods developed to aid understanding. The positive and negative approaches are not unlike cataphatic and apophatic Christian theology. In the first approach, imagery is used to describe the ineffable. The richness of such symbols is adequately described in several of Rene Guenon's works.

In the negative approach, Brahman is ``\emph{neti, neti}", or ``not this, not that". That is, no image or concept will ever be adequate. Once again, the ``head" will fail us, but the ``heart" will see. Griffiths reminds Christians of the same thing.

\paragraph{The Hidden Source}
Few people try to find the hidden source of the things in the phenomenal world. Once again, Griffiths illustrates this with a story about a Brahman youth who learned Sanskrit, memorized the Vedas, and so on. After his studies, the youth returned to his Father, who asked him:

\begin{quotex}
Svetaketu, as you are so conceited, considering yourself so well-read and so learned, my dear, have you ever learned of that instruction by which we hear that which cannot be heard, by which we perceive that which cannot be perceived, and by which we know that which we know that which cannot be known? 

\end{quotex}
The son did not know, so the father showed him a fruit from a banyan tree and opened it up. When asked, the son said he saw several seeds. When the father broke a seed and asked again, the son said he saw nothing. The father explained that the banyan tree arose from the subtle essence that cannot be seen.

An analogous idea in Eastern Christianity is that all things participate in God's creative energy through an inner principles, or \emph{logoi}. This is explained at \textit{The Uncreated Energy (logoi) of God in Nature}\footnote{\url{https://thoughtsintrusive.wordpress.com/2015/01/05/the-energy-of-god-in-nature}}. (H/T \textbf{Ekzy'l})

\paragraph{Thou Art That}
The deepest mystery of the Upanishads is ``\emph{Tat tvam asi}", ``Thou are That". This does not mean ``I am God", as some Western wannabe Vedantists have told me. That is how it must sound to the thinking rationalizing mind, but its mystical meaning is:

\begin{quotex}
I, in the deepest centre, the ground of my being, am one with that Brahman, the source of all creation. 

\end{quotex}
Although Griffiths does not express it in these terms, this shows us the two ways to the mystery of God. In the macrocosmic way, we look at the world and recognize God as the hidden source behind the world of phenomena. In the microcosmic way, we look within ourselves to find the source of our own being.

\paragraph{The Four States}
In following the path to the source of our own being, there are four states of consciousness. We identify the Self with one of those states, depending on what we understand to be real.

\subparagraph{Waking State}
In this state we take our ordinary waking state to be the most real. We are a body, we have experiences, we pursue money, sex and power, and so on. But all this passes away, so this cannot be the real Self.

\subparagraph{Dream State}
A man or woman may then try to find something more real, more immortal. So they look within. They learn to recognize their thoughts, feelings, desires, fears, likes, dislikes, and so on. At this point, people will be led to consider their Personality to be the real Self. Nevertheless, the personality will perish along with the body.

\subparagraph{Sleep State}
Those who go further will learn to detach from the personality, recognizing the real Self is neither the body nor the personality. At this point they may discover something deeper. They will see that much of their life is ``just happening". The body adopts certain postures and movements on its own. Thoughts and feelings arise from some deeper part of oneself, beyond conscious control. Hence the real Self is beyond any conscious awareness of it. This is the state of deep, or dreamless, sleep.

\subparagraph{Turiya}
The fourth state is beyond all the waking, dreaming, and sleep states, yet integrates them all. This is the true Self. We are so accustomed to believe that mind and reason are the supreme principles. But that is illusion (\emph{maya}) and ignorance (\emph{avidya}). Yet turiya is beyond thought.

\paragraph{Ego Death}
Griffiths relates the story of Nachiketas, who does down to the underworld and meets Yama, the god of Death. He explains:

\begin{quotex}
This is a fundamental principle of all religious teaching. If you want to reach your true Self, if you want to find God, you have to die. In the Christian tradition baptism is death. To be baptized is to participate in the death of Christ. The ego has to undergo this death, the ego, which is the person the mask which we seek to preserve. 

\end{quotex}
Nachiketas is offered three boons by Yama.

\begin{enumerate}
\item Nachiketas asked to be reconciled with his father. The seeker has left home in search of Truth. He eventually needs to return to be reconciled with his pat, with the tradition of his family and people. 
\item Nachiketas asks to understand the fir sacrifice, which is the offering of everything in the Cosmic fire of life. 
\item Nachiketas asks to know what lies beyond death. At first, Yama declines to answer, but finally explains about the Self, Atman, and so on. 
\end{enumerate}
Unfortunately, few are those today who would accept those boons. Words like ``God", ``soul" and so one have lost all significance as they don't correspond to experience. Even the religious, who may know doctrines, do not know God in their hearts.

\paragraph{The Abyss}
Real doctrine cannot be obtained by argument. Instead, Faith is necessary, when understood as an ``illumination of the mind". It comes by hearing. Beyond death by baptism, the illumination of faith, there is something else, as the Upanishads put it:

\begin{quotex}
The wise man, who, by means of meditation on his Self, recognizes the Ancient One who is difficult to be seen, who has entered into the dark, who is hidden in the cave, who dwells in the abyss as God; he indeed leaves joy and sorrow far behind. 

\end{quotex}
This is the experience of God in the darkness, and the Ancient One is the primeval source of your being.

\flrightit{Posted on 2015-03-08 by Cologero }

\begin{center}* * *\end{center}

\begin{footnotesize}\begin{sffamily}



\texttt{obscure on 2015-03-10 at 02:26 said: }

I notice you mention univocity, so I must warn you away from the narratives being produced by the `radical orthodoxy' movement wherein they derive some rubbish from Heidegger and Deleuze about Duns Scotus. I am referring of course to the `univocity of Being'. Their narratives however are not derived from familiarity with Scotus' work. Scotus argued that being is predicated univocally because the analogy is contained in the subject per se. Namely, Scotus argues that we begin metaphysics with an intuition of Being and this raises a disjunctive: finite or infinite? The mode of understanding the infinite subject is analogy and is not found in the mode of predication. Scotists call this doctrine `real analogy' as opposed to the older doctrine of `analogical predication'. However, scholars don't often read the fullness of the primary sources and are unfamiliar with whole systems; they are interested in narratives. Nothing new of course.

Scotus in a sense improves upon St. Thomas rather than contradicts him, hence he is called the `subtle doctor' while Thomas is the easier-to-grasp `common doctor'. Scotus also promoted the `absolute primacy of Christ'. which you may be interested in. The Franciscans promote it today, but Scotism is overshadowed by the more popular Neo-Thomism. Absolute primacy is a very important theological doctrine once it is understood adequately, since it entails that theosis is more fundamental than redemption (most important site for the doctrine):

\url{http://absoluteprimacyofchrist.org/}


\hfill

\texttt{Mark Citadel on 2015-03-10 at 23:43 said: }

It was through the stunning accuracy of the Doctrine of the Ages and the predictions of the Kali Yuga in the Vedic Tradition that led me to my current stand as a mild Christian Hermeticist (a less bold perennialism). There is simply no other way of explaining how this doctrine so accurately reflects the cyclical degeneration of mankind other than some higher revelation given to its author. Whether this was angelic or diabolical in nature is largely irrelevant, both angels and demons have their role in the order and balance of the Divine Realm, both being products of `primordial' origin. 

If one takes Christianity as the final revelation from God, he need not totally discount some elements of truth in other earlier faiths that may have had lower orders of divine influence, as long as they are not mutually exclusive. In fact, studying these faiths that also fed from the ever-present and almost riverine qualities of the World of Tradition, can lead to a greater appreciation and understanding of Scripture, of God, and of the Divine Realm. Even studying the perverse and fiendish cults of the Canaanites and such can give us insight into the nature of the demonic entities that presumed the status of gods over these debauched peoples and thus some knowledge of these entities that are beyond any realm of practical study and examination.

Great post by the way.


\hfill

\texttt{Saint-Efficace on 2015-03-19 at 06:41 said: }

``… the gods are names and forms under which the one God manifests himself. That is, they are more or less the equivalent of ``angels" or the ``cosmic powers" mentioned by St Paul."

Of course The Church does not believe and teach that angels are names and/or forms under which God manifests Himself, or any such nonsense. Quite the contrary. She also disagrees with Griffiths (as well as Guénon) on the inspired character of the Vedas. It is worth noting perhaps that the particular attempts at rapprochement of Catholic and Hindu positions of Griffiths c.s., however slight their resonance in the Catholic Church, have met with typical indifference and r a t h e r \ w o r s e among Hindus. Exceeding one's brief never was a good idea.

Anyway, Catholics, unlike (most) Hindus, still avoid like the plague any worship of angels as of false gods, ``quoniam omnes dii gentium daemonia" (Ps XCV: v) or, in the venerable English translation, ``For all the gods of the Gentiles are devils" (Douay-Rheims) following the Greek \textit{pantes hoi theoi to:n ethno:n daimonia} (LXX)

Like the modernisers you can have them as ``idols" if you want – which, if you will permit me, I venture to doubt you do, somehow – following the Hebrew: (kol elohei HaAmim elilim) of the trimmed MT. The reference is 96:5 in that case. Hardly an improvement for the case of the Devil's Party. Silver, gold, rocks, stones, devils – the Gentiles were happy enough to worship any of them. Your modern Gentlile, of course, desperately tries to join in, adding paper – but as often as not proves to be worshipping but the Self.


\hfill

\texttt{Cologero on 2015-03-29 at 00:48 said: }

It's my bad, Saint-Efficace. The question marks appeared because this blog does not support UTF8 encoding, which was not the default when it was created. It would take too much effort to change the encoding at this point.

So now we can move onto your bad. In your zeal to be super-correct, you have misread or misunderstood the text. Of course, Griffiths did not say that the angels are names or forms of God. That is the Vedic claim and he offered an alternative explanation, much as Thomas Aquinas did in respect to the Greeks. See Angels and Demons for an explanation.

Now I can quote the Bible as well as you, though probably not as well as the devil to which you refer. Following Griffiths, I offered a biblical explanation for a type of ``inspiration" and even pointed out that there is a greater revelation. That is perfectly orthodox teaching.

What Griffith does with the Vedic teaching is not unlike what the Fathers did with Plato and Thomas with Aristotle. Griffith's aims are not what you claim; in any case, the ideal of a ``rapprochement" was never mentioned in the post, so it is out-of-bounds in a cross-examination. Moreover, as we have explained previously, we deal with a text by (1) bringing out its logical conclusions and (2) placing it into a larger perspective. Please try to make an effort to see that and then reconsider your comment.

This larger perspective takes Augustine seriously when he claims there is but one tradition.


\hfill

Saint-Efficace on 2015-04-01 at 03:55 said:

Cologero, this being the Internet, I realise the mere text of a post may give more of an offensive impression than intended, so please, accept an aplogy. Not that I don’t stand by what I wrote; I merely intended it as a kind of reminder, not as an attack.
The devil, being the first ‘Plotestant’ as well as the first Jacobin can probably quote Scripture better than either of us. I have always agreed with St Augustine on the tradition.
A reading, time permitting, of an essay by Robert Fastiggi and Jose Pereira, titled “The Swami from Oxford” gives a good indication of where I tend to come from with respect to such issues and Fr Grifiths in particular after several decades of interest – summarising better than I could and saving me more typing than perhaps ought to be done in a comment box. It was in Crisis Magazine, but Catholic Culture has it here:

\url{http://www.catholicculture.org/culture/library/view.cfm?id=3427}

(I can’t help thinking you might also be interested in James Arraj’s work “Christianity in the Crucible of East-West Dialogue” which goes over much of the matter here discussed – and more besides. You may already know it, of course. Online at

\url{http://www.innerexplorations.com/catew/christia.htm})

Re St Augustine… As a reader of Tomberg you are no doubt aware of the interest in his work from the side of the Monastic Interreligious Dialogue/ Centering Prayer people; Fr Keating, Fr Pennington, Br Teasdale, Fr Barnhart… Some of the names, indeed, associated with Fr Griffiths and his legacy. Golden Cain newsletter, ‘Wisdom Christianity’? They appear to have it in for St Gus a bit. It’s early centuries yet, of course but I must admit to a measure of unease.

Anyway thanks for the reply.

\hfill


\texttt{argusandphoenix on 2015-03-17 at 22:41 said: }

Lewis has an interesting quote from Till We Have Faces:

``And in that far distant day when the gods become wholly beautiful, or we at last are shown how beautiful they always were, this will happen more and more (having dangerous enemies as one's own house). For mortals, as you said, will become more and more jealous. And mother and wife and child and friend will all be in league to keep a soul from being united with the Divine Nature…this age of ours will one day be the distant past. And the Divine Nature can change the past. Nothing is yet in its true Form."


\hfill

\texttt{Logres on 2015-03-18 at 08:22 said: }

This is his take on ``Grace Perfects Nature", and put in a poetic way. The density of the Kali Yuga is also a density of the gods mingling with men, which confuses things more, but has its own beauty.


\hfill

\texttt{Mark Citadel on 2015-04-17 at 13:33 said: }

Yet again, fascinating stuff.

``Brahman, Atman, and Purusha"

Remarkable how this number 3 appears over and over again across theologies, isn't it? `This Rough Beast' pointed it out just recently.


\hfill

\texttt{David on 2015-04-17 at 23:37 said: }

Those kind of text are most interesting when linked with other text such as "Christianity and the doctrine of non-duality". I highly suggest it. Great stuff again, thank you.


\hfill

\texttt{Cassiodorus on 2015-04-20 at 20:18 said: }

Christians, of course, affirm the reality of the soul. However, traditional Christianity also holds fast to the creature-Creator distinction and explicitly rejects the ``supreme identity". Is this merely a confessional bias or does this reflect a fundamental metaphysical first principle? Every time I think I have made peace with classical theism's relationship to the Perennial Philosophy, these concerns of mine come bubbling up to the surface yet again.


\hfill

\texttt{Cologero on 2015-04-20 at 20:24 said: }

Cassiodorus, Guenon claims that ``traditional Christianity" is a valid tradition in his sense. I suspect you are misunderstanding something. Specifically, it is not just a matter of verbal formulations but rather of a realization. You are over-thinking it.


\hfill

\texttt{Cassiodorus on 2015-04-20 at 22:19 said: }

Cologero,

As always, I appreciate your graciousness and patience in putting up with my repetitiveness. But, I'm troubled by the possibility that I'm not missing something. It seems to me that orthodox Christianity (contra the Perennialist School), aligns much more with Madhva's ``dualistic" school of Vedanta as opposed to Shankara's Advaita. 

I understand that esoterism transcends all discursivity (a contradiction perhaps). Nevertheless, the Christian mystical tradition distinguishes itself from the jnanic schools in that it speaks of participation, not identity. This is the case because of Christianty's strict adherence to the creature-Creator distinction. It seems to me that Palamite theology bears this out in stark relief. Even though the doctrine of divinization is affirmed (putting aside that this ,again, refers to union, not identity), the Eastern Church still has to distinguish between God's Essence and His Energies. The creature can join with His Energies, but NOT with His Essence, not with what He is in Himself. 

From my pov, the Supreme Identity, or ``realization" in a Shankaran sense, is simply not on the map for a classical theist. Perhaps, a case can be made for a qualified nondualism of a sort argued for by the likes of a Ramanuja. But I think that even this kind of panentheism might be too ``strong" to reconcile with Christian doctrine.

Kind Regards


\hfill

\texttt{David on 2015-04-21 at 04:15 said: }

Cassiodorius : Have you read Vladimir Lossky on the subject, or "Christianity and the doctrine of non-duality" by an anonymous cistercian monk ? After reading those, you will find that it is not totally alien to Christianity, simply said differently and under different terms. Is it the most prevalent view ? Of course not, but it is not in India either. Is it heterodox ? No, as explained in the cistercian monk book. And this work is much interesting because he said specifically that he doesn't want to delve into mysticism (such as Eckhart) because using this would make it more open to "heterodoxy" criticism, and so he goes to major thinker such as St Bernard, St Thomas Aquinas, St Albert, etc. You should also read Dyonisus the Aeropagite.


\hfill

\texttt{Cassiodorus on 2015-04-21 at 18:18 said: }

David, I appreciate the recommendations. I am familiar with those authors and works that you suggested. The book by the anonymous Trappist monk was a particularly interesting take on how one might look at the Incarnation from a Vedantin point of view. Nevertheless, it seems to me that is still very difficult to reconcile orthodox Christianity with Advaita Vedanta, that is, an unqualified nondualism. Perennialsts, such as Jean Borella, and Perennialist sympathizers such as Stratford Caldecott and Robert Bolton, have written on this subject. As I said in my previous comment, the Christian mystical streams never appears to stray from a fundamentally bhaktic, that is theistic, framework. But why?

" Christianity is surely the religion of creation ex nihilo and of the human person loved forever by a God who is different from himself." Moreover, ``Advaitins or perennialist non-dualists will find it hard to comprehend how the attainment of supreme realization might be made to depend on the incarnation of God at one particular point in history." – S. Caldecott

That being said, I don't think that a qualified nondualism poses the same problem. For example, Visistadvaita affirms that Reality is Brahman; however, individual souls are real and the universe is also real, being parts of Brahman or modes of His manifestation. This view, I think, ``saves the appearances"- can safeguard both the ``one and the many" and upholds our personhood as individuals.

\end{sffamily}\end{footnotesize}
